\section{API Design and Implementation}
\label{sec:api}

\subsection{RESTful Principles}

The API follows RESTful principles:

\begin{enumerate}
  \item \textbf{Stateless}: Each request contains all necessary information; the server stores no client context.
  \item \textbf{Cacheable}: Responses include appropriate caching directives.
  \item \textbf{Uniform Interface}: Resources are accessed via standard HTTP methods (GET, POST).
  \item \textbf{Resource Identification}: Endpoints represent resources (health, predictions), not actions.
  \item \textbf{Standard Representations}: JSON is the canonical representation format.
\end{enumerate}

\subsection{Endpoints}

\subsubsection{Health Check Endpoint}

\begin{verbatim}
  GET /health
\end{verbatim}

\textbf{Purpose:} Report the operational status of the API and model.

\textbf{Response (200 OK):}
\begin{verbatim}
  {
    "status": "healthy" | "unhealthy",
    "model_loaded": true | false
  }
\end{verbatim}

\textbf{Semantics:}
\begin{itemize}
  \item \code{status}: "healthy" if the model is loaded and ready, "unhealthy" otherwise.
  \item \code{model\_loaded}: Boolean indicating whether the model artifact was successfully loaded.
\end{itemize}

\textbf{Use Cases:}
\begin{itemize}
  \item Kubernetes liveness/readiness probes.
  \item Load balancer health checks.
  \item Monitoring and alerting systems.
  \item Manual inspection during troubleshooting.
\end{itemize}

\subsubsection{Single Prediction Endpoint}

\begin{verbatim}
  POST /predict
\end{verbatim}

\textbf{Request (Content-Type: application/json):}
\begin{verbatim}
  {
    "user_id": "196",
    "movie_id": "242"
  }
\end{verbatim}

\textbf{Response (200 OK):}
\begin{verbatim}
  {
    "user_id": "196",
    "movie_id": "242",
    "predicted_rating": 4.12,
    "model_version": "1.0.0"
  }
\end{verbatim}

\textbf{Error Responses:}
\begin{itemize}
  \item \textbf{422 Unprocessable Entity}: Request validation failed (missing fields, invalid types).
  \item \textbf{503 Service Unavailable}: Model not loaded or unavailable.
  \item \textbf{500 Internal Server Error}: Unexpected error during prediction (e.g., model corruption).
\end{itemize}

\textbf{Semantics:}
\begin{itemize}
  \item \code{user\_id}, \code{movie\_id}: Echo the request inputs (enables request tracing).
  \item \code{predicted\_rating}: Floating-point prediction in $[1.0, 5.0]$.
  \item \code{model\_version}: Version identifier of the model used (from environment config).
\end{itemize}

\subsubsection{Batch Prediction Endpoint (Bonus)}

\begin{verbatim}
  POST /predict/batch
\end{verbatim}

\textbf{Request:}
\begin{verbatim}
  {
    "predictions": [
      {"user_id": "196", "movie_id": "242"},
      {"user_id": "196", "movie_id": "300"},
      ...
    ]
  }
\end{verbatim}

\textbf{Response (200 OK):}
\begin{verbatim}
  {
    "predictions": [
      {"user_id": "196", "movie_id": "242", 
       "predicted_rating": 4.12, "model_version": "1.0.0"},
      {"user_id": "196", "movie_id": "300", 
       "predicted_rating": 3.87, "model_version": "1.0.0"},
      ...
    ],
    "total_count": N
  }
\end{verbatim}

\textbf{Rationale:} Batch predictions reduce HTTP overhead and enable efficient client-side caching strategies. 
The endpoint returns results in the same order as requests, ensuring deterministic client handling.

\subsubsection{Information Endpoints}

\begin{itemize}
  \item \code{GET /}: Root endpoint returning API metadata (name, version, docs URL).
  \item \code{GET /model/info}: Model metadata (version, type, loaded status).
  \item \code{GET /docs}: Swagger UI (auto-generated by FastAPI).
  \item \code{GET /redoc}: ReDoc UI (alternative documentation format).
  \item \code{GET /openapi.json}: OpenAPI schema in JSON.
\end{itemize}

These endpoints are standard in modern APIs and are automatically provided by FastAPI.

\subsection{Request/Response Schemas}

Request and response validation is implemented using Pydantic:

\subsubsection{PredictionRequest Schema}

\begin{verbatim}
  class PredictionRequest(BaseModel):
      user_id: str = Field(..., min_length=1, 
                           examples=["196"],
                           description="User ID")
      movie_id: str = Field(..., min_length=1, 
                            examples=["242"],
                            description="Movie ID")
\end{verbatim}

\textbf{Validation:}
\begin{itemize}
  \item \code{user\_id} and \code{movie\_id} are required non-empty strings.
  \item Missing fields result in 422 error.
  \item Type mismatches (e.g., numeric instead of string) are rejected.
  \item Examples are used in automatic documentation.
\end{itemize}

\subsubsection{PredictionResponse Schema}

\begin{verbatim}
  class PredictionResponse(BaseModel):
      user_id: str
      movie_id: str
      predicted_rating: float = Field(..., ge=1.0, le=5.0,
                                      description="Predicted rating (1.0-5.0)")
      model_version: str
\end{verbatim}

\textbf{Validation:}
\begin{itemize}
  \item \code{predicted\_rating} is constrained to $[1.0, 5.0]$ via Pydantic validators.
  \item Serialization to JSON enforces types.
\end{itemize}

\subsubsection{HealthResponse Schema}

\begin{verbatim}
  class HealthResponse(BaseModel):
      status: str
      model_loaded: bool
\end{verbatim}

\subsection{Error Handling}

The API implements comprehensive error handling:

\begin{verbatim}
  @app.post("/predict")
  async def predict(request: PredictionRequest):
      if model is None or not model.is_loaded():
          raise HTTPException(status_code=503, 
                              detail="Model not loaded")
      try:
          rating = model.predict(request.user_id, 
                                 request.movie_id)
          return PredictionResponse(...)
      except Exception as e:
          logger.exception(f"Prediction error: {e}")
          raise HTTPException(status_code=500, 
                              detail=str(e))
\end{verbatim}

\textbf{Error Semantics:}
\begin{itemize}
  \item \textbf{422}: Validation errors (invalid request structure).
  \item \textbf{503}: Service degradation (model unavailable).
  \item \textbf{500}: Unexpected errors (model corruption, system issues).
\end{itemize}

All errors return JSON with a \code{detail} field, enabling client-side error handling.

\subsection{CORS and Middleware}

The API enables Cross-Origin Resource Sharing (CORS) for browser-based clients:

\begin{verbatim}
  app.add_middleware(
      CORSMiddleware,
      allow_origins=["*"],
      allow_credentials=True,
      allow_methods=["*"],
      allow_headers=["*"],
  )
\end{verbatim}

\textbf{Production Note:} The wildcard configuration (\code{allow\_origins=["*"]}) is convenient for 
development but permissive. Production deployments should restrict origins to known clients.

\subsection{API Documentation}

FastAPI automatically generates interactive API documentation:

\begin{itemize}
  \item \textbf{Swagger UI} (\code{/docs}): Interactive API explorer with try-it-out capability.
  \item \textbf{ReDoc} (\code{/redoc}): Alternative documentation format emphasizing readability.
  \item \textbf{OpenAPI Schema} (\code{/openapi.json}): Machine-readable API specification.
\end{itemize}

These are automatically derived from:
\begin{itemize}
  \item Docstrings in endpoint handlers.
  \item Type hints and Pydantic schemas.
  \item Field descriptions and examples.
\end{itemize}

This reduces documentation drift: the documentation always matches the implementation.
